%% capitoli/cap1_inquadramento.tex
%% In testa al file, prima di scrivere, incollare l'estratto di
%% output/F4-indice.md relativo a questo capitolo (tesi sostenuta,
%% fonti, lunghezza stimata): serve a controllare gli scostamenti.
%% Convenzione delle etichette: sez:cap1:<nome>, tab:cap1:<nome>.

% Regola di proporzione: questo capitolo non supera il 15%% del lavoro.
% E' premessa, non oggetto.

\chapter{La funzione consultiva della Corte dei conti nel sistema costituzionale}
\label{cap:cap1}

\section{L'art. 100 Cost. e la doppia natura della Corte}
\label{sez:cap1:art100}


\section{La radice storica: il parere al Governo nel t.u. del 1934}
\label{sez:cap1:storia}


\section{Ausiliarietà, neutralità, collaborazione}
\label{sez:cap1:fondamento}


\section{La contabilità pubblica come materia: perimetro e confini mobili}
\label{sez:cap1:materia}


